\documentclass{beamer}
\usepackage{tikz,amsmath,amssymb,hyperref,graphicx,stackrel,animate,tipa}
\usetikzlibrary{positioning,shadows,arrows,shapes,calc,dsp,chains}
\newcommand{\argmax}{\operatornamewithlimits{argmax}}
\newcommand{\argmin}{\operatornamewithlimits{argmin}}
\mode<presentation>{\usetheme{Frankfurt}}
\AtBeginSection[]
{
  \begin{frame}<beamer>
    \frametitle{Outline}
    \tableofcontents[currentsection,currentsubsection]
  \end{frame}
}
\title{Lecture 26: Linear Prediction}
\author{Mark Hasegawa-Johnson}
\date{ECE 401: Signal Processing}  
\begin{document}

% Title
\begin{frame}
  \maketitle
\end{frame}

% Title
\begin{frame}
  \tableofcontents
\end{frame}


%%%%%%%%%%%%%%%%%%%%%%%%%%%%%%%%%%%%%%%%%%%%%%%%%%%%%%%%%%%%%%%%%%%%%
\section[LPC]{Linear prediction}

\begin{frame}
  \frametitle{Review: Covariance}

  \begin{itemize}
  \item The \textbf{covariance} between samples $x[m]$ and $x[m-n]$ is
    \begin{align*}
      \rho_x[m,m-n]  = E\left[(x[m]-\mu_x[m])(x[m-n]-\mu_x[m-n])^*\right]
    \end{align*}
  \item Covariance is interesting because it tells us the best linear prediction
    of sample $x[m]$ given sample $x[m-n]$:
    \begin{align*}
      \hat{x}[m] &= \mu_x[m] + \frac{\rho_x[m,m-n]}{\sigma^2_x[m-n]}
      \left(x[m-n]-\mu_x[m-n]\right)
    \end{align*}
  \end{itemize}
\end{frame}

\begin{frame}
  \frametitle{Review: Autocorrelation approximates covariance}

  \begin{itemize}
  \item Autocorrelation:
  \begin{align*}
    \frac{r_x[n]}{N-n} =\frac{1}{N-n}\sum_{m=n}^{N-1} x[m]x^*[m-n] \approx E\left[x[m]x^*[m-n]\right]
  \end{align*}
  \item Covariance:
    \begin{align*}
      \rho_x[m,m-n]  = E\left[(x[m]-\mu_x[m])(x[m-n]-\mu_x[m-n])^*\right]
    \end{align*}
  \item For zero-mean, stationary signals, covariance is the expected average value of
    autocorrelation
  \end{itemize}
\end{frame}


\begin{frame}
  \frametitle{Linear predictive coding (LPC): Synthesis model}

  A linearly predictable signal is one whose samples, $s[n]$, can be estimated
  as the linear prediction of $p$ previous samples plus some small error term $e[n]$:
  \begin{align*}
    s[n] = e[n]+\sum_{k=1}^p a_k s[n-k]
  \end{align*}
\end{frame}

\begin{frame}
  \frametitle{LPC prediction model}

  We can say that $\tilde{s}[n]$ is the predictable part of the signal:
  \begin{align*}
    \tilde{s}[n] &= \sum_{k=1}^p a_k s[n-k]
  \end{align*}
  \ldots and the whole signal is the sum of its predictable part plus the error.
  For this reason the error is sometimes called the excitation:
  \begin{align*}
    s[n] &= \tilde{s}[n]+e[n]\\
  \end{align*}
\end{frame}

\begin{frame}
  \frametitle{The LPC analysis filter}

  Another way to say the same thing: We model the error, $e[n]$,
  as being the part of $s[n]$ that cannot be predicted based on a
  linear combination of the previous $p$ samples:
  \begin{align*}
    e[n] &= s[n]-\tilde{s}[n]\\
    &= s[n]-\sum_{k=1}^p a_k s[n-k]
  \end{align*}
\end{frame}

%%%%%%%%%%%%%%%%%%%%%%%%%%%%%%%%%%%%%%%%%%%%%%%%%%%%%%%%%%%%%%%%%%%%%
\section[Coefficients]{Solving for the LPC Coefficients: General Form}

\begin{frame}
  \frametitle{Mean-squared error (MSE)}

  We will choose $a_k$ to minimize the mean squared error (MSE):
  \begin{align*}
    \mathcal{E} &= \sum_m e^2[m]\\
    &= \sum_m \left(s[m]-\sum_{k=1}^p a_ks[m-k]\right)^2
  \end{align*}
\end{frame}

\begin{frame}
  \frametitle{Minimize the MSE}

  \begin{align*}
    \mathcal{E} &= \sum_m \left(s[m]-\sum_{k=1}^p a_ks[m-k]\right)^2
  \end{align*}
  Differentiate $\mathcal{E}$ with respect to $a_i$:
  \begin{align*}
    \frac{\partial \mathcal{E}}{\partial a_i} 
    &= 2\sum_m s[m-i]\left(s[m]-\sum_{k=1}^p a_ks[m-k]\right)
  \end{align*}
  \ldots and set the derivative to zero:
  \begin{align*}
    0 &= 2\sum_m s[m-i]\left(s[m]-\sum_{k=1}^p a_ks[m-k]\right)
  \end{align*}
  We get the solution:
  \begin{align*}
    \sum_m s[m-i]a[m] &= \sum_{k=1}^p a_k \sum_m s[m-k]s[m-i]
  \end{align*}
\end{frame}

\begin{frame}
  \frametitle{LPC: The Solution!}

  So we discover that, to minimize $\mathcal{E}=\sum_me^2[m]$, we
  should choose the coefficients $a_i$ so that
  \begin{align*}
    \sum_m s[m-i]s[m] &= \sum_{k=1}^p a_k\sum_m s[m-k]s[m-i]
  \end{align*}
  Notice that each of those summations is an autocorrelation!  The
  optimum LPC coefficients are given by
  \begin{align*}
    r_s[i] &= \sum_{k=1}^p a_kr_s[k-i]
  \end{align*}
  (notice that $r_s[k-i]=r_s[i-k]$)
\end{frame}

\begin{frame}
  \frametitle{The LPC solution as a matrix equation}

  So we discover that, to minimize $\mathcal{E}=\sum_me^2[m]$, we
  should choose the coefficients $a_i$ so that
  \begin{align*}
    r_s[i] &= \sum_{k=1}^p a_kr_s[k-i]
  \end{align*}
  Now, notice that this is $p$ linear equations (one for each $i$) in
  $p$ unknowns ($a_1$ through $a_p$), so we can write it as a single
  matrix equation:
  \begin{align*}
    \boldsymbol{\gamma} &= \mathbf{R}\mathbf{a}
  \end{align*}
  \begin{align*}
    \boldsymbol{\gamma}=\left[\begin{array}{c}r_s[1]\\\vdots\\r_s[p]\end{array}\right],~~~
    \mathbf{R}=\left[\begin{array}{ccc}r_s[0]&\cdots&r_s[p-1]\\
        \vdots&\ddots&\vdots\\r_s[p-1]&\cdots&r_s[0]\end{array}\right],~~~
    \mathbf{a}=\left[\begin{array}{c}a_1\\\vdots\\a_p\end{array}\right]
  \end{align*}
\end{frame}
    
\begin{frame}
  \frametitle{The LPC solution as a matrix equation}

  So we discover that, to minimize $\mathcal{E}=\sum_me^2[m]$, we
  only need to invert a matrix!
  \begin{align*}
    \mathbf{a}=\mathbf{R}^{-1}\boldsymbol{\gamma}
  \end{align*}
  \begin{align*}
    \mathbf{a}=\left[\begin{array}{c}a_1\\\vdots\\a_p\end{array}\right],~~~
    \mathbf{R}=\left[\begin{array}{ccc}r_s[0]&\cdots&r_s[p-1]\\
        \vdots&\ddots&\vdots\\r_s[p-1]&\cdots&r_s[0]\end{array}\right],~~~
    \boldsymbol{\gamma}=\left[\begin{array}{c}r_s[1]\\\vdots\\r_s[p]\end{array}\right],~~~
  \end{align*}
\end{frame}
    
\begin{frame}
  \frametitle{Try the quiz!}

  Go to the course web page, and try the quiz!
\end{frame}

\begin{frame}
  \frametitle{What is the residual error $\mathcal{E}$?}

  Now we might ask: What is $\mathcal{E}$ after we've
  minimized it?
  \begin{align*}
    \mathcal{E} &= \sum_me^2[m]\\
    &= \sum_m \left(s[m]-\sum_{k=1}^p a_k s[m-k]\right)^2\\
    &= r_s[0] - 2\sum_k a_kr_s[k] + \sum_i\sum_k a_ia_kr_s[k-i]
  \end{align*}
  If we plug in $r_s[k]=\sum_ia_ir_s[k-i]$, it simplifies to
  \begin{align*}
    \mathcal{E} &= r_s[0] - \sum_k a_kr_s[k]
  \end{align*}
\end{frame}

\begin{frame}
  \frametitle{What is the residual error $\mathcal{E}$?}

  Notice that $r_s[0]$ is the energy of the speech signal \textbf{before} prediction:
  \begin{align*}
    r_s[0] = \sum_m s^2[m]
  \end{align*}
  So we can say that $\sum_ka_kr_s[k]$ is the portion of the original energy that
  is explained by linear prediction:
  \begin{align*}
    \mathcal{E} &= r_s[0] - \sum_k a_kr_s[k] = r_s[0]-\mathbf{a}^T\boldsymbol{\gamma}\\
    \mathbf{a}=\left[\begin{array}{c}a_1\\\vdots\\a_p\end{array}\right],~~~
    &\boldsymbol{\gamma}=\left[\begin{array}{c}r_s[1]\\\vdots\\r_s[p]\end{array}\right]
  \end{align*}
    
\end{frame}

%%%%%%%%%%%%%%%%%%%%%%%%%%%%%%%%%%%%%%%%%%%%%%%%%%%%%%%%%%%%%%%%%%%%%
\section[Implementations]{Implementations: Autocovariance and Autocorrelation Methods}

\begin{frame}
  \frametitle{$p$ linear equations in $p$ unknowns}

  Remember that the general solution is:
  \begin{align*}
    \mathbf{a} &=\mathbf{R}^{-1}\boldsymbol{\gamma}
  \end{align*}
  where $\mathbf{R}$ is a $p\times p$ matrix.  Computing its inverse
  is $\mathcal{O}\left\{p^3\right\}$ (I've read that solutions exist
  that are $\mathcal{O}\left\{p^{2.3}\right\}$, but normal Gaussian
  elimination is $\mathcal{O}\left\{p^3\right\}$).  We might ask: Is
  there any way to reduce the computation?
\end{frame}
  
\begin{frame}
  \frametitle{Autocorrelation $\Rightarrow$ Toeplitz matrix}

  If we window the signal before computing the autocorrelation, then
  $r_s[i-k]$ only depends on $|i-k|$, so the matrix $\mathbf{R}$ has a
  special structure called \textbf{Toeplitz}:
  \begin{align*}
    \boldsymbol{\gamma} &= \mathbf{R}\mathbf{a}
  \end{align*}
  \begin{align*}
    \mathbf{R}=\left[\begin{array}{cccc}
        r_s[0]&r_s[1]&\cdots&r_s[p-1]\\
        r_s[1]&r_s[0]&\cdots&r_s[p-2]\\
        \vdots&\vdots&\ddots&\vdots\\
        r_s[p-1]&r_s[p-2]&\cdots&r_s[0]\end{array}\right]
  \end{align*}
  An equation $\mathbf{a}=\mathbf{R}^{-1}\boldsymbol{\gamma}$ with a Toeplitz
  $\mathbf{R}$ can be solved in $\mathcal{O}\left\{p^2\right\}$ using
  a special equation called the Levinson-Durbin recursion.  
\end{frame}

\begin{frame}
  \frametitle{What are the disadvantages of the autocorrelation solution?}

  The disadvantage of autocorrelation is that we are predicting each
  sample only using previous samples of the \textbf{windowed} speech
  signal:
  \begin{align*}
    e[m] &= s[m]-\sum_k a_ks[m-k]
  \end{align*}
  \begin{itemize}
  \item The sample at the beginning of the window, $s[0]$, is being
    predicted based on samples outside the window, which are zero: $s[-p]=s[-1]=0$!
  \item Since the error has this built-in artifact, the coefficients
    will have some built-in error.
  \item This problem is most important if $w[n]$ is short.
  \end{itemize}
\end{frame}

\begin{frame}
  \frametitle{The autocovariance solution: Predict from samples outside the window}

  \begin{itemize}
  \item  There is another solution, used more often in financial predictions (and less
    often in signal processing) called the ``autocovariance solution.'' 
    The autocovariance solution is to predict each sample possibly using
    samples from before the start time of the window, i.e., use some
    known nonzero value for $s[-1]$ and $s[-p]$.
  \item Autocovariance gives more accurate coefficients for short windows, but\ldots
  \item The matrix is no longer Toeplitz (e.g.,
    $\phi_n[1,1]\ne\phi_n[2,2]$), so solving
    $\mathbf{a}=\mathbf{R}^{-1}\boldsymbol{\gamma}$ is an
    $\mathcal{O}\left\{p^3\right\}$ operation.
  \item More important in real-world applications: autocorrelation
    gives a synthesis filter that is guaranteed stable, but autocovariance doesn't.
  \item For those reasons, we'll only use autocorrelation in this
    class, not autocovariance.
  \end{itemize}
\end{frame}


%%%%%%%%%%%%%%%%%%%%%%%%%%%%%%%%%%%%%%%%%%%%%%%%%%%%%%%%%%%%%%%%%%%%%%%%%%5
\section{LPC-10 Speech Synthesis Model}


\begin{frame}
  \frametitle{Speech synthesis using LPC}

  \begin{itemize}
  \item LPC-10 was a secure (easily encryptable) speech coding standard used in the 1980s
  \item It used an impulse train (for voiced sounds) or white noise
    (for unvoiced sounds) as the input to a 10-tap LPC prediction filter
  \item You can see the flowchart and hear examples at \url{https://nmikstas.github.io/portfolio/talkingClock/talkingClock.html}
  \end{itemize}
\end{frame}


%%%%%%%%%%%%%%%%%%%%%%%%%%%%%%%%%%%%%%%%%%%%%%%%%%%%%%%%%%%%%%%%%%%%%%%%%%5
\section{Summary}

\begin{frame}
  \frametitle{Summary}

  \begin{itemize}
  \item LPC models the resonant structure of speech using a
    $p^{\text{th}}$-order linear predictor:
    \begin{align*}
      s[n] = e[n]+\sum_{k=1}^p a_k s[n-k]
    \end{align*}
  \item Minimizing the mean-squared error results in $p$ linear
    equations in $p$ unknowns, which can be solved as
    \begin{align*}
      \boldsymbol{\gamma}=\mathbf{R}\mathbf{a}
    \end{align*}
    \begin{align*}
      \boldsymbol{\gamma}=\left[\begin{array}{c}r_s[1]\\\vdots\\r_s[p]\end{array}\right],~~
      \mathbf{R}=\left[\begin{array}{ccc}r_s[0]&\cdots&r_s[p-1]\\
          \vdots&\ddots&\vdots\\r_s[p-1]&\cdots&r_s[0]\end{array}\right],~~
      \mathbf{a}=\left[\begin{array}{c}a_1\\\vdots\\a_p\end{array}\right]
    \end{align*}
  \end{itemize}
\end{frame}

\end{document}

